\documentclass{article}
\usepackage[T1]{fontenc}
\usepackage{lmodern}
\usepackage{graphicx}
\usepackage{amsmath,amssymb}
\usepackage[a4paper,margin=2cm]{geometry}
\usepackage[absolute,overlay]{textpos}
\setlength{\TPHorizModule}{1mm}
\setlength{\TPVertModule}{1mm}
\usepackage[indonesian]{babel}
\usepackage{hyperref}
\setlength{\emergencystretch}{2em}
% unit: Halaman 52

\begin{document}

% === Halaman 52 (python/native) ===

52

• Pada Contoh 2,


- kunci = abcdef 


- panjang kunci = 6


- jarak antara dua crypto yang berulang = 16


- 16 bukan kelipatan 6

 crypto tidak dienkripsi menjadi kriptogram yang sama

• Goal metode Kasiski: mencari dua atau lebih kriptogram yang 

berulang untuk menentukan panjang kunci.

Plainteks   : cryptoisshortforcryptography

Kunci  

   : abcdefabcdefabcdefabcdefabcd

Cipherteks: CSASXTITUKWSTGQUCWYQVRKWAQJB

16

\end{document}
